% One-time imported draft. Editable native source; no Markdown regeneration.
\section{Step 23 Explain the integer error
  relation}\label{007-bresenham-line-s23-step-23-explain-the-integer-error-relation}

Let i be x displacement from the ordered start and j be signed y
displacement along sy. The state relation is err = dx \ensuremath{-} D \ensuremath{-} D\ensuremath{\times}i + dx\ensuremath{\times}j.
An x step subtracts D; a y step adds dx. The fixed relation connects the
update instructions to the line geometry.

\subsection{Write or
  inspect}\label{007-bresenham-line-s23-write-or-inspect}

Use the derivation fixture to explain one x-only update and one diagonal
update.

\subsection{Expected
  observation}\label{007-bresenham-line-s23-expected-observation}

Both coordinate decisions use one frozen e2, then update a state
consistent with the new coordinates.

\subsection{Check your
  understanding}\label{007-bresenham-line-s23-check-your-understanding}

Why is err not simply the perpendicular distance in pixels?

\begin{itemize}
  \tightlist
  \item
        \textbf{A} It is stored in a register.
  \item
        \textbf{B} It is a scaled decision expression with an initial offset
        and signed contributions.
  \item
        \textbf{C} It can only hold addresses.
\end{itemize}
